\section*{Bab \ref{ch:b1:chemical-biogenic-sediments} --- Sedimen karbonat, silika dan evaporit}

\begin{solution}{exo:b1:chemical-biogenic-sediments:1}
Klorida $0.5503\times35.165 = \qty{19.35}{g}$; natrium $0.3066\times35.165 =
\qty{10.78}{g}$; sulfat \qty{2.711}{g}; magnesium \qty{1.284}{g}. Natrium dan
klorida bersama-sama membentuk $0.5503 + 0.3066 = 0.857$, yakni \qty{85.7}{\%} dari
garam.
\end{solution}

\begin{solution}{exo:b1:chemical-biogenic-sediments:2}
Fraksi yang tersisa adalah $1/\mathrm{CF}$: kalsium karbonat $1/1.8 = 0.556$ tersisa,
\qty{44.4}{\%} menguap; gipsum $0.263$ tersisa, \qty{73.7}{\%} menguap; halit $0.094$
tersisa, \qty{90.6}{\%} menguap; magnesium sulfat $0.014$ tersisa, \qty{98.6}{\%} menguap.
\end{solution}

\begin{solution}{exo:b1:chemical-biogenic-sediments:3}
(a)~$\rho_s(1-\phi) = G/a = 8.4/0.0042 = \qty{2000}{kg.m^{-3}}$. (b)~Dengan
aragonit, $1 - 2000/2930 = 0.32$, porositas \qty{32}{\%}; dengan kalsit,
$1 - 2000/2710 = 0.26$. (c)~$a = 3.7/2000 = \qty{1.85e-3}{m}$, \qty{1.85}{mm} setahun.
\end{solution}

\begin{solution}{exo:b1:chemical-biogenic-sediments:4}
Hingga \qty{4000}{m} kandungan turun hanya perlahan (88 hingga 80); di bawah
\qty{4000}{m} ia turun tajam: \omterm{def:b1:chemical-biogenic-sediments:ccd}{lisoklin} terletak sekitar \qty{4000}{m}. Kandungan turun
melewati \qty{10}{\%} di antara \qty{4700}{m} (12) dan \qty{4800}{m} (6):
$4700 + 100\times(12-10)/(12-6) = \qty{4733}{m}$, CCD sekitar \qty{4730}{m}, sekitar
\qty{700}{m} di bawah \omterm{def:b1:chemical-biogenic-sediments:ccd}{lisoklin}.
\end{solution}

\begin{solution}{exo:b1:chemical-biogenic-sediments:5}
(a)~\ce{Ca^2+ + 2 HCO3- -> CaCO3 + CO2 + H2O}: satu mol \ce{CO2} per mol
\ce{CaCO3}, sehingga $44.01/100.09 = \qty{0.440}{kg}$ \ce{CO2} per kilogram
\ce{CaCO3}. (b)~$3.7\times0.440 = \qty{1.63}{kg}$ \ce{CO2} per meter persegi dan per
tahun. (c)~Dari dua ion hidrogen karbonat yang dipakai, hanya satu karbon yang masuk
ke \omterm{def:b1:rock-forming-minerals:mineral}{mineral}; yang lain keluar sebagai karbon dioksida. Kalsifikasi mengambil alkalinitas
dari air dan menurunkan kapasitasnya untuk menampung karbon terlarut; berapa banyak
\ce{CO2} yang sampai ke udara bergantung pada penyangga air laut, suatu soal kesetimbangan
karbonat dalam buku-buku kimia.
\end{solution}

\begin{solution}{exo:b1:chemical-biogenic-sediments:6}
(a)~$t_c = (2000/350)^2 = \qty{32.7}{Ma}$; pada \qty{25}{km} per juta tahun,
$25\times32.7 = \qty{816}{km}$ dari sumbu. (b)~$t_c = (2600/350)^2 =
\qty{55.2}{Ma}$, \qty{1380}{km}. (c)~Dasar berumur \qty{60}{Ma} telah menerima
lempung merah selama $60 - 32.7 = \qty{27.3}{Ma}$ pada kasus pertama, dan hanya
selama $60 - 55.2 = \qty{4.8}{Ma}$ pada kasus kedua.
\end{solution}

\begin{solution}{exo:b1:chemical-biogenic-sediments:7}
(a)~Ya: $a = 4.2 > r = 3$, dan \omterm{def:b1:chemical-biogenic-sediments:reef}{terumbu} mendapat \qty{1.2}{mm} setahun atas laut hingga
puncaknya mencapai permukaan. (b)~$D$ tumbuh sebesar $10 - 4.2 = \qty{5.8}{mm}$ setahun;
dari 2 hingga \qty{50}{m} membutuhkan $48/0.0058 = \num{8300}$ tahun. (c)~Karena laju
pertumbuhannya turun ketika air makin dalam: lebih sedikit cahaya mencapai alga yang
hidup di dalam karang dan memberi mereka makan, sehingga $a$ berkurang dan \omterm{def:b1:chemical-biogenic-sediments:reef}{terumbu}
tertinggal makin cepat.
\end{solution}

\begin{solution}{exo:b1:chemical-biogenic-sediments:8}
(a)~Reaksi yang terjadi adalah \ce{CaSO4.2H2O -> CaSO4 + 2 H2O}. (b)~Volume molarnya adalah
$172.17/2.30 =
\qty{74.9}{cm^3}$ dan $136.14/2.97 = \qty{45.8}{cm^3}$, sehingga padatan kehilangan
$1 - 45.8/74.9 = \qty{38.8}{\%}$ volumenya. (c)~$2\times18.02/172.17 = 0.209$:
\qty{20.9}{\%} massa keluar sebagai air. Dilepaskan di kedalaman, air ini menaikkan
tekanan dalam pori-pori batuan di sekitarnya; di bawah penutup ia dapat membuat batuan
lebih lemah dan membantu sesar tergelincir, dan ia dapat melarutkan serta membawa garam lain.
\end{solution}

\begin{solution}{exo:b1:chemical-biogenic-sediments:9}
(a)~$9.2/255 = \qty{3.6}{\%}$ terawetkan, \qty{96.4}{\%} larut. (b)~Di tempat produksi
sangat tinggi, beberapa persen darinya tetap merupakan fluks yang besar, dan di tempat
hampir tidak ada lempung atau karbonat yang mengencerkannya: di bawah air naik di
Samudra Selatan, Pasifik utara dan Pasifik khatulistiwa, dan di bawah CCD. (c)~Frustul
yang tipis dan rapuh larut lebih dulu dan yang kokoh bertahan: sedimen menyimpan seleksi
spesies yang paling kokoh, bukan komunitas yang hidup di air.
\end{solution}

\begin{solution}{exo:b1:chemical-biogenic-sediments:10}
(a)~$EA = 1.5\times\qty{1e10}{m^2} = \qty{1.5e10}{m^3}$ setahun. (b)~Dengan
$Q_{\mathrm{in}} = \qty{1.8e10}{m^3}$, $Q_{\mathrm{out}} = \qty{0.3e10}{m^3}$ dan
$k = 6$: melewati ambang gipsum (3.8) tetapi tidak ambang halit (10.6), sehingga kalsium
karbonat dan gipsum mengendap. Dengan $Q_{\mathrm{in}} =
\qty{1.575e10}{m^3}$, $Q_{\mathrm{out}} = \qty{0.075e10}{m^3}$ dan $k = 21$, untuk ion
yang tetap dalam larutan seperti magnesium: air asin melewati ambang halit, sehingga
halit juga mengendap, tetapi bukan garam magnesium (70). (c)~Halit membutuhkan
$Q_{\mathrm{in}}/(Q_{\mathrm{in}} - EA) \ge 10.6$, yakni $Q_{\mathrm{in}} \le 10.6\,EA/9.6 = 1.104\,EA = \qty{1.66e10}{m^3}$ setahun: aliran masuk boleh melebihi
penguapan paling banyak sekitar sepersepuluh.
\end{solution}

\begin{solution}{exo:b1:chemical-biogenic-sediments:11}
(a)~Air laut mengandung \qty{67.3}{mg} bromida per kilogram; pada 10.6~kali, air asin
mengandung $67.3\times10.6 = \qty{713}{mg.kg^{-1}}$ dan halit pertama
$0.032\times713 = \qty{22.8}{mg.kg^{-1}}$. (b)~Pada 30~kali, $67.3\times30 =
\qty{2019}{mg.kg^{-1}}$ dalam air asin dan \qty{64.6}{mg.kg^{-1}} dalam halit.
(c)~Halit hanya mengambil 3.2 perseratus dari konsentrasi air asin, sehingga hampir seluruh
bromida tetap dalam larutan dan terus terkonsentrasi ketika air menguap. (d)~Kenaikan
mencatat konsentrasi air asin; penurunan mendadak mencatat pengenceran olehnya, oleh
masuknya air laut baru ke dalam cekungan, atau oleh air tawar, sebelum halit terus tumbuh.
\end{solution}

\begin{solution}{exo:b1:chemical-biogenic-sediments:12}
(a)~$t_c = \qty{32.65}{Ma}$, sehingga $0.02\times\num{32650} = \qty{653}{m}$ ooze
(\num{32650} interval seribu tahun pada \qty{2}{cm}) dan
$0.001\times\num{37350} = \qty{37}{m}$ lempung merah. (b)~$t_c = (2500/350)^2 =
\qty{51.0}{Ma}$: \qty{1020}{m} ooze dan \qty{19}{m} lempung merah. CCD yang lebih
dalam menghasilkan ooze yang jauh lebih tebal, dan lempung merah dasar Atlantik yang
dalam lebih tipis daripada Pasifik dengan umur yang sama. (c)~CCD mungkin sempat lebih
dalam untuk sementara, sehingga situs itu naik kembali di atasnya; atau lempeng membawa
situs di bawah sabuk produktivitas tinggi, seperti khatulistiwa, tempat lebih banyak
karbonat berjatuhan; atau lapisan itu dibawa oleh longsoran atau \omterm{def:b1:clastic-sediments:turbidite}{arus turbidit} dari
bagian dasar laut yang lebih tinggi.
\end{solution}

\begin{solution}{pb:b1:chemical-biogenic-sediments:1}
\textbf{1.} $1030\times1000 = \qty{1.03e6}{kg}$ air laut, yang mengandung
$1.03\times10^6\times0.035165 = \qty{36200}{kg}$ garam.

\textbf{2.} Gipsum $1.03\times10^6\times1.581\times10^{-3} = \qty{1628}{kg}$;
halit $1.03\times10^6\times0.02741 = \qty{28230}{kg}$.

\textbf{3.} Kalsium karbonat $113.3/2710 = \qty{0.042}{m}$; gipsum $1628/2300 =
\qty{0.71}{m}$; halit $28230/2170 = \qty{13.0}{m}$.

\textbf{4.} Anhidrit memiliki massa $1628\times136.14/172.17 =
\qty{1288}{kg}$ dan ketebalan $1288/2970 = \qty{0.43}{m}$.

\textbf{5.} Kering, garamnya adalah $0.110 + 1.250 + 27.41 + 6.34 = \qty{35.11}{g}$ per
kilogram, dari jumlah itu halit adalah $27.41/35.11 = \qty{78}{\%}$: natrium dan klorida
membentuk sebagian besar garam terlarut, dan halit mengambil hampir semuanya.

\textbf{6.} Dari bawah ke atas: kalsium karbonat, gipsum, halit, garam pahit. Setiap
\omterm{def:b1:rock-forming-minerals:mineral}{mineral} mulai mengendap ketika air asin mencapai kejenuhannya, yang paling sukar larut
lebih dulu, dan mengendap di atas apa yang diendapkan sebelumnya.

\textbf{7.} Kedalaman air asin adalah $1000/\mathrm{CF}$: \qty{556}{m} untuk kalsium
karbonat, \qty{263}{m} untuk gipsum, \qty{94}{m} untuk halit.

\textbf{8.} $1 - 1/10.6 = \qty{90.6}{\%}$.

\textbf{9.} $1000/70 = \qty{14}{m}$ untuk magnesium sulfat, $1000/90 = \qty{11}{m}$
untuk garam kalium.

\textbf{10.} Sekitar \qty{1.10}{g.cm^{-3}} ketika kalsium karbonat terbentuk, 1.13 dengan
gipsum, 1.22 dengan halit dan lebih dari 1.30 dengan garam pahit, dibandingkan dengan
1.03 untuk air laut. Air laut yang mengalir masuk mengapung di atas air asin yang padat,
yang tetap di dasar; dalam cekungan berpenghalang air asin dapat mengalir kembali ke
samudra di bawah aliran masuk.

\textbf{11.} Lapisan itu hanya \qty{4}{cm} tebalnya, dibandingkan dengan \qty{13}{m}
halit; dan ia terbentuk di tepi, tempat air yang masuk pertama kali terkonsentrasi,
lebih daripada di tengah, sehingga ia tercampur dengan dasar gipsum atau digantikan olehnya.

\textbf{12.} Halit hanya mengambil sebagian kecil bromida air asin (latihan~\ref{exo:b1:chemical-biogenic-sediments:11}), sehingga air asin menjadi makin kaya bromida seiring halit tumbuh, dan setiap lapisan halit baru mengambil lebih banyak.

\textbf{13.} $1500/13.01 = 115.3$: 116 pengisian.

\textbf{14.} $116\times0.434 = \qty{50}{m}$ anhidrit, bukan sebagai satu lapisan tetapi
sebagai 116 lapisan, masing-masing sedikit kurang dari setengah meter tebalnya, di dasar
halit dari pengisiannya: suksesi siklik anhidrit dan halit.

\textbf{15.} Per pengisian, $1.03\times10^6\times6.34\times10^{-3}/2000 =
\qty{3.3}{m}$, sehingga \qty{379}{m} untuk 116 pengisian. Garam-garam ini yang paling
mudah larut: aliran masuk air laut berikutnya melarutkannya kembali, dan cekungan yang
airnya keluar sebelum akhir tidak pernah mencapainya.

\textbf{16.} Pada setiap pengisian ulang air asin diencerkan: kandungan bromida turun
tajam di dasar halit baru dan naik lagi ke atas, sebuah profil gigi gergaji dengan satu
gigi untuk setiap pengisian.

\textbf{17.} $2000/13.01 = 154$ pengisian.

\textbf{18.} $5000\times365.25\times86400 = \qty{1.578e11}{m^3}$ setahun, dibagi di atas
$\qty{1.0e11}{m^2}$: $E = \qty{1.58}{m}$ setahun.

\textbf{19.} $1.578\times10^{11}\times1030\times0.035165 = \qty{5.72e12}{kg}$ garam setahun.

\textbf{20.} $1.578\times10^{11}\times1030\times0.02741/2170 =
\qty{2.05e9}{m^3}$ halit setahun, yakni \qty{0.0205}{m}, sekitar \qty{2}{cm} setahun di
atas dasar.

\textbf{21.} $1500/0.0205 = \num{73000}$ tahun.

\textbf{22.} Permukaan turun sebesar \qty{1.58}{m} setahun; \qty{1000}{m} air hilang
dalam $1000/1.578 = 634$ tahun.

\textbf{23.} Halit muncul ketika \qty{90.6}{\%} air telah hilang, setelah
$905.7/1.578 = 574$ tahun.

\textbf{24.} \qty{1.5}{km} halit membutuhkan garam dari 116 pengisian cekungan (115.3,
dibulatkan ke atas), baik jika mereka datang satu demi satu saat cekungan mengering dan
terisi kembali, maupun terus-menerus melalui ambang, seperti dalam soal 21, pada
laju satu cekungan air laut setiap $1000/1.578 = 634$ tahun.
\end{solution}
